\section{Critical exponents at a computable transition}\label{sec:intro}
Local interactions with rational coefficients can have an uncomputable
asymptotic exponent at a rational critical coupling. The organizing idea
is to make the coupling an energetic ruler. A classical layer selects
cells of pitch $p_*(g)=2r/g+O(1)$; a quantum counter inside each cell
then makes the thermodynamic gap equal to a finite box's gap. Finite-size
scaling becomes an exact spectral identity in a deliberately constructed
model. The size of the local alphabet is fixed for each prescribed
exponent and does not grow as $g$ approaches the transition.

Critical gap exponents at this computable transition fill $\mathsf D_2$
in the bulk on $\Z^2$ and $\mathsf D_3$ in the bulk on $\Z^3$,
as well as at a surface-anchored film. The two-dimensional attainment
has a computably tamed volume limit. In three dimensions an internally
created frame supports an untamed classical computation; a final-time
broadcast sends its settled word to quantum counters throughout the
volume. The distinction is in how the finite counter labels are selected,
not in whether the gap-carrying modes occupy the bulk. These are
attainment statements in the gap convention below, not a dimension-only
ceiling for every two-dimensional family.

We use the gap above the \emph{entire} ground space, namely the difference
between the two lowest distinct eigenvalues. A scalar finite matrix has
no such gap and is excluded whenever a gap is required. A family has
translation-invariant onsite and oriented nearest-neighbour terms
$h_a(g)=h_a^0+g h_a^1$, with rational Hermitian matrices. Only the onsite
term needs to depend on $g$ in our constructions. At rational $g$ every
finite-volume matrix is algebraic. On tori, take the commensurate sizes
$L=n!$; on the wetting open cubes we also take $L=n!$; the fixed-thickness film uses all $L\in\N$. Define, in this order,
\begin{equation}\label{eq:critical-definition}
 \begin{gathered}
 \Delta_\infty(g)=\lim_{L\to\infty}\Delta(L,g),\qquad
 x(g)=\frac{\log(1/\Delta_\infty(g))}{\log(1/g)},\\
 (z\nu)^-=\liminf_{g\downarrow0}x(g),\qquad
 (z\nu)^+=\limsup_{g\downarrow0}x(g).
 \end{gathered}
\end{equation}
The first limit must exist and be positive at each sufficiently small
positive coupling. An existing finite exponent is $z\nu=(z\nu)^-=(z\nu)^+$.
General arithmetic ceilings concern rational couplings. In the explicit
rulers the gap is constant between rational transition points, so real
and rational approaches have the same endpoints. The notation $z\nu$
denotes this gap exponent; it does not assert a dispersion relation or
identify a ground-state correlation length.

For an individual real, write $\mathsf S_m$ and $\mathsf P_m$ for the
left-$\Sigma^0_m$ and left-$\Pi^0_m$ classes of its strict lower rational
cut. Write $\mathsf D_m=\Delta^0_m$ for reals computable in
$0^{(m-1)}$. Rational endpoints are treated separately when passing
between strict upper and lower cuts. Uniform families throughout the
paper are specified by oracle Cauchy names or one-sided approximation
names, as explained in Section~\ref{sec:arithmetic}; uniform decidability
of exact comparison with a rational is not assumed.

\begin{theorem}[Bulk ruler]\label{thm:bulk-headline}
For every $0\leq a\leq b<\infty$ with $a\in\mathsf S_2$ and
$b\in\mathsf P_2$, a fixed rational nearest-neighbour affine family on
$\Z^2$ has $((z\nu)^-,(z\nu)^+)=(a,b)$ along commensurate tori.
The transition coupling is $g_c=0$. The gap is at least $1/16$ for
$g\leq0$, and tends to zero from the right. For small positive $g$,
the ground configurations fill the plane with cells and
\[
 \Delta_\infty(g)=D(p_*(g)-1),\qquad p_*(g)=\frac{2r}{g}+O(1),
 \qquad r=\frac1{16},
\]
off the explicitly specified tie points. Here $D(h)$ is the full quantum
counter gap selected by an $h\times h$ box. At ties the gap is the
minimum of the two neighbouring values. The ground energy density is
$0$ for $g\leq0$ and $-g^2/(4r)+O(g^4)$ for $g>0$.
The inner gap limit is computably tamed. Consequently its finite
nonnegative existing exponents are exactly $\mathsf D_2$.
\end{theorem}

The proof is in Sections~\ref{sec:ruler} and \ref{sec:transfer}, using
the finite-size construction in Sections~\ref{sec:statement}--\ref{sec:necessity}.
In particular $z\nu$ can be the halting probability $\Omega_U$ of any
chosen prefix-universal machine. No uncomputable number is an interaction
coefficient. There is one counter in every cell, with a positive density
of possible excitation locations at fixed $g$. These gap-closing modes
are not supported on the boundary of a finite box. We make no
identification with a boundary-independent operator-algebraic gap.


\begin{theorem}[Three-dimensional bulk broadcast]\label{thm:bulk-three}
For every finite $\alpha\geq0$ in $\mathsf D_3$, a fixed
translation-invariant nearest-neighbour rational affine family on $\Z^3$
has critical gap exponent $\alpha$ on cubic tori $L=n!$, with the volume
limit taken first and the entire ground space removed. The limiting gap
is zero for $g\leq0$, positive for $0<g<11/480$, and closes as
$g\downarrow0$. The ground energy density is zero for $g\leq0$ and
$-g^2/(4r)+O(g^4)$ from the right, with $r=1/16$.
At each fixed positive coupling in that interval, the gap is carried by
at least $L^3/p^2$ independent bulk reader excitations for all sufficiently
large factorial $L$, where $p=2r/g+O(1)$ is a gap-attaining pitch.
The corresponding local spectral weight per volume is at least $1/p^2$.
The finite-volume gap is eventually exact at fixed $g$, without a
general computable settling bound. Together with the general two-limit
ceiling, the existing finite nonnegative exponents in this convention
fill exactly $\mathsf D_3$. A selected infinite-volume phase has the
same gap above its full zero-energy subspace.
\end{theorem}

Section~\ref{sec:bulk-three} gives the complete construction and proof,
including the minimum-of-two-gaps rule at pitch ties. The internally
anchored computation does not confine the soft quantum modes to a seam:
their possible locations have positive density at fixed $g$. No unique
ground state, all-integer-size limit or equality with a locally
nondegenerate or KMS-extremized gap is asserted.

\begin{theorem}[Wetting ruler and film hierarchy]\label{thm:film-headline}
A surface-anchored rational affine family along factorial open cubes in $\Z^3$
selects a wetting-film height $h_*(g)=(r/g)^{1/2}+O(1)$.
Its finite nonnegative existing gap exponents fill $\mathsf D_3$, and
its ordered endpoint pairs fill $\mathsf S_3\times\mathsf P_3$.
The gapless side includes $g=0$. At fixed $g>0$ the finite-box gap is
eventually exactly constant, without a general computable settling bound.
Independently, for integer interactions the fixed-thickness film
exponents, with the in-plane limit taken first, fill exactly the same
classes. More generally $k$ dependent nested box limits, including the
outer scaling limit, fill $\mathsf D_{k+1}$ and its ordered
$\mathsf S_{k+1}\times\mathsf P_{k+1}$ pairs in dimension $k+1$.
\end{theorem}

The wetting statement and its precise selection rules are in
Section~\ref{sec:wetting}; the separate fixed-thickness theorem and the
completed finite local-rule schema are in Sections~\ref{sec:film} and
\ref{sec:nested}. Necessity requires every intermediate limit to exist
for every remaining finite parameter tuple, and only asserts a class
when the relevant final endpoints are finite. Their surface selection remains distinct from the bulk broadcast
of Theorem~\ref{thm:bulk-three}.

The examples are finely tuned. For the earlier two-dimensional and
wetting rulers, a specified arbitrarily small onsite perturbation
changes the programmed gap exponent to zero
(Section~\ref{sec:instability}). The ruler pitch has
structural exponent one; a theorem making a conventional correlation
length exponent $\nu$ uncomputable is not claimed. Finite analytic
systems have integer gap exponents. This forces the effective finite
system to change infinitely often when its envelope has a noninteger
exponent, but does not prove that all noninteger critical exponents in
physics require a staircase.

\subsection{Relation to earlier work}
Spectral-gap undecidability for translation-invariant lattices was
established by Cubitt, P\'erez-Garc\'ia and Wolf
\cite{spectralShort,spectralFull}. Bausch, Cubitt and Watson constructed
uncomputable phase diagrams \cite{phaseDiagrams}. Purcell, Li and Cubitt
placed a transition at an uncomputable coupling \cite{chaitinTransition}.
Here the coupling is rational and an asymptotic gap exponent contains
the uncomputable number. Watson, Onorati and Cubitt proved that a
computable renormalisation map can have uncomputably complex flows
\cite{renormalisation}. None of these statements identifies stability
with an effective convergence modulus.

Fisher and Barber developed finite-size scaling for critical phenomena
\cite{fisherBarber}. Our cell decomposition supplies a particular exact
transfer, rather than a theorem about the usual scaling ansatz in all
models. The commensurate steps belong to the established staircase
tradition of Fisher and Selke and of Bak \cite{fisherSelke,bak}; the
arithmetic programming of their gap envelope is the additional feature.
Rellich's analytic perturbation theory, as presented by Kato
\cite{kato}, supplies the finite-system obstruction.

The open-edge preference is a form of the boundary-bonus mechanism of
Bausch, Cubitt, Lucia, P\'erez-Garc\'ia and Wolf \cite{sizeDriven}; the
exact flag identities and all-sector bounds used here are proved below.
Watson and Cubitt study complexity of ground energy density
\cite{energyDensity}. Xu et al. distinguish intrinsic bulk excitations
from boundary-dependent gaps \cite{bulkGap}. Our bulk broadcast
has positive-density local excitations, but its gap removes all
zero-energy states. Section~\ref{subsec:bulk-conventions} separates
this convention from their locally nondegenerate, KMS-extremized gap. The real hierarchy and density
classifications are prior mathematical results of Zheng and Weihrauch
and of Downey, Jockusch and Schupp \cite{zhengWeihrauch,density}.
Meyerovitch classified entropy-dimension values of multidimensional
shifts of finite type using arithmetical classes of reals
\cite{meyerovitchGrowth}. This is a prior classification of local-rule
growth exponents; the observable and realization here concern spectral
gaps and their coupling-selected finite-cell envelopes.
The companion precision problem \cite{pricePrecision} is independent
of the present proof.

There are closer precedents than undecidable transition locations.
Schmidhuber's algorithmic ensemble has the computable critical
inverse temperature $\log2$ and partition value $\Omega_U$ there
\cite{algorithmicTransition}; it is not a local quantum-lattice
gap-exponent classification. Balister, Bollob\'as, Morris and Smith
explicitly announce uncomputability of bootstrap critical-probability
exponents through a halting reduction \cite{bootstrapUniversality}.
That announcement, referring to work in preparation, concerns failure
of a uniform exponent algorithm across update families; it does not
by itself exhibit one fixed noncomputable real quantum gap exponent.
Watson and Bausch study the complexity of approximating critical
boundaries under a finite-size consistency promise
\cite{criticalBoundary}, a distinct input-size complexity question.
Castilla-Castellano and Lucia show that arbitrarily small onsite
perturbations remove undecidable behavior from a particular
one-dimensional gap family \cite{gapInstability}. Our explicit
perturbation proves the quantitative change of a programmed exponent
to zero; neither result proves computability of every stable exponent.
No exhaustive priority claim is made.
